% !TEX root = main.tex
After a short introduction to the supervised learning of convnets, we describe our unsupervised approach as well as the specificities of its optimization.

\subsection{Preliminaries}
Modern approaches to computer vision, based on statistical learning, require good image featurization.
In this context, convnets are a popular choice for mapping raw images to a vector space of fixed dimensionality.
When trained on enough data, they constantly achieve the best performance on standard classification benchmarks~\cite{he2015delving,krizhevsky2012imagenet}.
We denote by $f_\theta$ the convnet mapping, where $\theta$ is the set of corresponding parameters.
We refer to the vector obtained by applying this mapping to an image as feature or representation.
Given a training set $X = \{x_1, x_2, \dots , x_N\}$ of $N$ images, we want to find a parameter $\theta^*$ such that the mapping $f_{\theta^*}$ produces good general-purpose features.

These parameters are traditionally learned with supervision, \ie each image $x_n$ is associated with a label $y_n$ in $\{0,1\}^k$.
This label represents the image's membership to one of $k$ possible predefined classes.
A parametrized classifier $g_W$ predicts the correct labels on top of the features $f_\theta(x_n)$.
The parameters $W$ of the classifier and the parameter $\theta$ of the mapping are then jointly learned by optimizing the following problem:
\begin{eqnarray}\label{eq:sup}
  \min_{\theta, W} \frac{1}{N} \sum_{n=1}^N \ell\left(g_W\left( f_\theta(x_n) \right), y_n\right),
\end{eqnarray}
where $\ell$ is the multinomial logistic loss, also known as the negative log-softmax function.
This cost function is minimized using mini-batch stochastic gradient descent~\cite{bottou2012stochastic} and backpropagation to compute the gradient~\cite{lecun1998gradient}.


\subsection{Unsupervised learning by clustering}
When $\theta$ is sampled from a Gaussian distribution, without any learning, $f_\theta$ does not produce good features.
However the performance of such random features on standard transfer tasks, is far above the chance level.
For example, a multilayer perceptron classifier on top of the last convolutional layer of a random AlexNet achieves 12\% in accuracy on ImageNet while the chance is at $0.1\%$~\cite{noroozi2016unsupervised}.
The good performance of random convnets is intimately tied to their convolutional structure which gives a strong prior on the input signal.
The idea of this work is to exploit this weak signal to bootstrap the discriminative power of a convnet.
We cluster the output of the convnet and use the subsequent cluster assignments as ``pseudo-labels'' to optimize Eq.~(\ref{eq:sup}).
This deep clustering (DeepCluster) approach iteratively learns the features and groups them.

Clustering has been widely studied and many approaches have been developed for a variety of circumstances. In the absence of points of comparisons,
we focus on a standard clustering algorithm, $k$-means. Preliminary results with other clustering algorithms indicates that this choice is not crucial.
$k$-means takes a set of vectors as input, in our case the features $f_\theta(x_n)$ produced by the convnet, and clusters them into $k$ distinct groups based on a geometric criterion.
More precisely, it jointly learns a $d\times k$ centroid matrix $C$ and the cluster assignments $y_n$ of each image $n$ by solving the following problem:
\begin{equation}
\label{eq:kmeans}
  \min_{C \in \mathbb{R}^{d\times k}}
  \frac{1}{N}
  \sum_{n=1}^N
  \min_{y_n \in \{0,1\}^{k}}
  \| f_\theta(x_n) -  C y_n \|_2^2
  \quad
  \text{such that}
  \quad
  y_n^\top 1_k = 1.
\end{equation}
Solving this problem provides a set of optimal assignments $(y_n^*)_{n\le N}$ and a centroid matrix $C^*$.
These assignments are then used as pseudo-labels; we make no use of the centroid matrix.

Overall, \OURS alternates between clustering the features to produce pseudo-labels using Eq.~(\ref{eq:kmeans})
and updating the parameters of the convnet by predicting these pseudo-labels using Eq.~(\ref{eq:sup}).
This type of alternating procedure is prone to trivial solutions; we describe how to avoid such degenerate solutions in the next section.


\subsection{Avoiding trivial solutions}
The existence of trivial solutions is not specific to the unsupervised training of neural networks, but to any method that jointly learns a discriminative classifier and the labels.
Discriminative clustering suffers from this issue even when applied to linear models~\cite{xu2005maximum}.
Solutions are typically based on constraining or penalizing the minimal number of points per cluster~\cite{bach2008diffrac,joulin2012convex}.
These terms are computed over the whole dataset, which is not applicable to the training of convnets on large scale datasets.
In this section, we briefly describe the causes of these trivial solutions and give simple and scalable workarounds.
\\

\noindent\textbf{Empty clusters.}
A discriminative model learns decision boundaries between classes.
An optimal decision boundary is to assign all of the inputs to a single cluster~\cite{xu2005maximum}.
This issue is caused by the absence of mechanisms to prevent from empty clusters and arises in linear models as much as in convnets.
A common trick used in feature quantization~\cite{johnson2017billion} consists in automatically reassigning empty clusters during the $k$-means optimization.
More precisely, when a cluster becomes empty, we randomly select a non-empty cluster and use its centroid with a small random perturbation as the new centroid for the empty cluster.
We then reassign the points belonging to the non-empty cluster to the two resulting clusters.
\\

\noindent\textbf{Trivial parametrization.}
If the vast majority of images is assigned to a few clusters, the parameters $\theta$ will exclusively discriminate between them.
In the most dramatic scenario where all but one cluster are singleton, minimizing Eq.~(\ref{eq:sup}) leads to a
trivial parametrization where the convnet will predict the same output regardless of the input.
This issue also arises in supervised classification when the number of images per class is highly unbalanced.
For example, metadata, like hashtags, exhibits a Zipf distribution, with a few labels dominating the whole distribution~\cite{joulin2016learning}.
A strategy to circumvent this issue is to sample images based on a uniform distribution over the classes, or pseudo-labels.
This is equivalent to weight the contribution of an input to the loss function in Eq.~(\ref{eq:sup}) by the inverse of the size of its assigned cluster.


\subsection{Implementation details}

\noindent\textbf{Convnet architectures.}
For comparison with previous works, we use a standard AlexNet~\cite{krizhevsky2012imagenet} architecture.
It consists of five convolutional layers with $96$, $256$, $384$, $384$ and $256$ filters; and of three fully connected layers.
We remove the Local Response Normalization layers and use batch normalization~\cite{ioffe2015batch}.
We also consider a VGG-16~\cite{simonyan2014very} architecture with batch normalization.
Unsupervised methods often do not work directly on color and different strategies have been considered as alternatives
~\cite{doersch2015unsupervised,noroozi2016unsupervised}.
We apply a fixed linear transformation based on Sobel filters to remove color and increase local contrast~\cite{bojanowski2017unsupervised,paulin2015local}.
\\

\noindent\textbf{Training data.}
We train \OURS on ImageNet~\cite{deng2009imagenet} unless mentioned otherwise.
It contains~$1.3$M images uniformly distributed into~$1,000$ classes.
\\

\noindent\textbf{Optimization.}
We cluster the central cropped images features and perform data augmentation (random horizontal flips and crops of random sizes and aspect ratios) when training the network. 
This enforces invariance to data augmentation which is useful for feature learning~\cite{dosovitskiy2014discriminative}.
The network is trained with dropout~\cite{srivastava2014dropout}, a constant step size, an~$\ell_2$
penalization of the weights~$\theta$ and a momentum of~$0.9$.
Each mini-batch contains $256$ images.
For the clustering, features are PCA-reduced to $256$ dimensions, whitened and $\ell_2$-normalized.
We use the $k$-means implementation of Johnson~\etal~\cite{johnson2017billion}.
Note that running k-means takes a third of the time because a forward pass on the full dataset is needed.
One could reassign the clusters every $n$ epochs, but we found out that our setup on ImageNet (updating the clustering every epoch) was nearly optimal.
On Flickr, the concept of epoch disappears: choosing the tradeoff between the parameter updates and the cluster reassignments is more subtle.
We thus kept almost the same setup as in ImageNet.
We train the models for $500$ epochs, which takes $12$ days on a Pascal P100 GPU for AlexNet.
\\

\noindent\textbf{Hyperparameter selection.}
We select hyperparameters on a down-stream task,~\ie, object classification on
the validation set of \textsc{Pascal} VOC with no fine-tuning.
We use the publicly available code of Kr\"ahenb\"uhl\footnote{https://github.com/philkr/voc-classification}.

